\section{BIJLAGE · Spiekbrief en resultaten}

\begin{missie}
Missie: alles op één blad. Gebruik het tijdens het oefenen, leer het voor de toets.
\end{missie}

\subsection*{Stappen van een onderzoek}

1 probleemstelling – 2 onderzoeksvraag en hypothese – 3 onderzoek opzetten en uitvoeren – 4 data verzamelen – 5 data analyseren en samenvatten – 6 conclusies trekken – 7 antwoord op de onderzoeksvraag – 8 reflecteren en communiceren.

\subsection*{Variabelen en data}

\breedtes{32/34/34}
\begin{tabular}{lll}
\hline
Variabele & Data & Voorbeeld \\
\hline
kwantitatief & numeriek & leeftijd, aantal uren \\
kwalitatief & categorisch geordend & tevredenheid \\
kwalitatief & categorisch niet-geordend & postcode, afdeling \\
\hline
\end{tabular}

\subsection*{Steekproef en populatie}

Populatie ($N$) = alle elementen waarover je iets wil weten. Steekproef ($n$) = deel dat effectief wordt onderzocht. Redenen voor een steekproef: te duur of te veel tijd, ethisch, onderzoek vernietigt de populatie. Steekproefvariabiliteit: verschillende steekproeven geven verschillende resultaten. Verzamelende – beschrijvende – verklarende statistiek.

\subsection*{Representatief?}

Drie criteria: voldoende groot – samenstelling komt overeen met de populatie – niet te veel non-respons. Aselect = door toeval; EAS = elk element dezelfde kans. Vertekend door: opportunistische steekproef, vrijwillige respons, non-respons.

\subsection*{Misleidingen}

Grafiek: as niet bij 0, 3D-diagram, beelddiagram. Cijfers: gemiddelde of mediaan, procent of procentpunt: $\dfrac{\text{nieuw}-\text{oud}}{\text{oud}}\cdot 100\,\%$. Data verzamelen: vraagstelling, opdrachtgever, representativiteit.

\subsection*{Samenhang en causaliteit}

Samenhang $\neq$ causaliteit. Misvattingen: omgekeerde oorzakelijke richting, toeval, derde variabele. Observationeel (niet ingrijpen, enkel samenhang) of experimenteel (ingrijpen, randomisatie, experimentele groep en controlegroep: causaliteit mogelijk).

\subsection*{Mijn resultaten}

\breedtes{14/46/20/20}
\begin{tabular}{llcc}
\hline
Level & Onderwerp & Score CHECK & Behaald? \\
\hline
1 & Statistisch onderzoek & \_\_ / 8 & ☐ \\
2 & Steekproef en populatie & \_\_ / 9 & ☐ \\
3 & Steekproefontwerp & \_\_ / 9 & ☐ \\
4 & Misleidingen en valkuilen & \_\_ / 8 & ☐ \\
5 & Samenhang versus causaliteit & \_\_ / 9 & ☐ \\
6 & Eindmissie & & ☐ \\
Eind & Eindtoets & \_\_ / 60 & ☐ \\
\hline
\end{tabular}
